\FigureLayoutDeclare{img/2001/new_curriculum_liberal/q11_fig1.png}{9.5cm}{0bp 0bp 4bp 0bp}

\chapter{2001年新课标卷（文）}
\section{单选题}
\begin{problem}
设 \(A = \{x \mid x^2 - x = 0\}\)，\(B = \{x \mid x^2 + x = 0\}\)，则 \(A \cap B\) 等于（\quad）

\choices
    {\(0\)}
    {\(\{0\}\)}
    {\(\varnothing\)}
    {\(\{-1,0,1\}\)}
\end{problem}

\begin{answer}
B
\end{answer}

\begin{solution}
由 \(x^2-x=x(x-1)\)，得 \(A=\{0,1\}\)；由 \(x^2+x=x(x+1)\)，得 \(B=\{-1,0\}\). 因此 \(A\cap B=\{0\}\)，故选 B.
\end{solution}

\begin{problem}
若 \(S_n\) 是数列 \(\{a_n\}\) 的前 \(n\) 项和，且 \(S_n = n^2\)，则 \(\{a_n\}\) 是（\quad）

\choices
    {等比数列，但不是等差数列}
    {等差数列，但不是等比数列}
    {等差数列，而且也是等比数列}
    {既非等比数列又非等差数列}
\end{problem}

\begin{answer}
B
\end{answer}

\begin{solution}
由 \(a_1=S_1=1\)，以及 \(a_n=S_n-S_{n-1}\)（\(n\geq2\)），得
\(a_n=n^2-(n-1)^2=2n-1\)（\(n\geq2\)）. 这个式子对 \(n=1\) 也成立，所以 \(a_n=2n-1\). 因而 \(a_{n+1}-a_n=2\)，数列是等差数列. 若它还是等比数列，则前三项应满足 \(a_2^2=a_1a_3\)，但 \(a_2^2=9\neq5=a_1a_3\)，故不是等比数列，选 B.
\end{solution}

\begin{problem}
过点 \(A(1,-1)\)，\(B(-1,1)\) 且圆心在直线 \(x + y - 2 = 0\) 上的圆的方程是（\quad）

\choices
    {\((x - 3)^2 + (y + 1)^2 = 4\)}
    {\((x + 3)^2 + (y - 1)^2 = 4\)}
    {\((x - 1)^2 + (y - 1)^2 = 4\)}
    {\((x + 1)^2 + (y + 1)^2 = 4\)}
\end{problem}

\begin{answer}
C
\end{answer}

\begin{solution}
线段 \(AB\) 的中点为 \(M(0,0)\)，直线 \(AB\) 的斜率为 \(-1\)，所以其垂直平分线为 \(y=x\). 圆心既在 \(y=x\) 上，又在 \(x+y-2=0\) 上，解得圆心为 \(C(1,1)\). 半径平方为
\(CA^2=(1-1)^2+(1+1)^2=4\). 因此圆的方程为 \((x-1)^2+(y-1)^2=4\)，选 C.
\end{solution}

\begin{problem}
若定义在区间 \((-1,0)\) 内的函数 \(f(x)=\log_{2a}(x+1)\) 满足 \(f(x)>0\)，则 \(a\) 的取值范围是（\quad）

\choices
    {\(\left(0, \frac{1}{2}\right)\)}
    {\(\left(0, \frac{1}{2}\right]\)}
    {\(\left(\frac{1}{2}, +\infty\right)\)}
    {\((0, +\infty)\)}
\end{problem}

\begin{answer}
A
\end{answer}

\begin{solution}
当 \(-1<x<0\) 时，\(0<x+1<1\). 对任意 \(t\in(0,1)\)，有 \(\log_c t>0\) 当且仅当 \(0<c<1\). 因此必须有 \(0<2a<1\)，即 \(0<a<\frac12\). 这也保证对数的底数满足 \(2a>0\) 且 \(2a\neq1\)，故选 A.
\end{solution}

\begin{problem}
若向量 \(\bs{a}=(3,2)\)，\(\bs{b}=(0,-1)\)，则向量 \(2\bs{b}-\bs{a}\) 的坐标是（\quad）

\choices
    {\((3,-4)\)}
    {\((-3,4)\)}
    {\((3,4)\)}
    {\((-3,-4)\)}
\end{problem}

\begin{answer}
D
\end{answer}

\begin{solution}
\(2\bs{b}=(0,-2)\)，所以
\(2\bs{b}-\bs{a}=(0,-2)-(3,2)=(-3,-4)\). 故选 D.
\end{solution}

\begin{problem}
设 \(A\)，\(B\) 是 \(x\) 轴上的两点，点 \(P\) 的横坐标为 \(2\) 且 \(\left|PA\right|=\left|PB\right|\). 若直线 \(PA\) 的方程为 \(x-y+1=0\)，则直线 \(PB\) 的方程是（\quad）

\choices
    {\(x + y - 5 = 0\)}
    {\(2x - y - 1 = 0\)}
    {\(2y - x - 4 = 0\)}
    {\(2x + y - 7 = 0\)}
\end{problem}

\begin{answer}
A
\end{answer}

\begin{solution}
由 \(P\) 的横坐标为 \(2\) 且 \(P\) 在直线 \(x-y+1=0\) 上，得 \(P=(2,3)\). 直线 \(PA\) 与 \(x\) 轴交于 \(A=(-1,0)\)，从而
\(PA^2=(2+1)^2+3^2=18\). 设 \(B=(t,0)\)，由 \(PB=PA\) 得
\((t-2)^2+3^2=18\)，即 \((t-2)^2=9\)，所以 \(t=-1\) 或 \(t=5\). 因为 \(A\)、\(B\) 是两点，故 \(B=(5,0)\). 过 \(P(2,3)\)、\(B(5,0)\) 的直线为 \(x+y-5=0\)，选 A.
\end{solution}

\begin{problem}
若 \(0 < \alpha < \beta < \frac{\pi}{4}\)，\(\sin \alpha + \cos \alpha = a\)，\(\sin \beta + \cos \beta = b\)，则（\quad）

\choices
    {\(a < b\)}
    {\(a > b\)}
    {\(ab < 1\)}
    {\(ab > 2\)}
\end{problem}

\begin{answer}
A
\end{answer}

\begin{solution}
设 \(g(t)=\sin t+\cos t\). 当 \(0<t<\frac\pi4\) 时，
\(g'(t)=\cos t-\sin t>0\)，所以 \(g(t)\) 在该区间上单调递增. 由 \(\alpha<\beta\)，得到
\(a=g(\alpha)<g(\beta)=b\)，故选 A.
\end{solution}

\begin{problem}
若椭圆经过原点，且焦点为 \(F_1(1,0)\)，\(F_2(3,0)\)，则其离心率为（\quad）

\choices
    {\(\dfrac{3}{4}\)}
    {\(\dfrac{2}{3}\)}
    {\(\dfrac{1}{2}\)}
    {\(\dfrac{1}{4}\)}
\end{problem}

\begin{answer}
C
\end{answer}

\begin{solution}
椭圆中心为两焦点的中点 \((2,0)\)，焦半距为 \(c=1\). 原点在椭圆上，且
\(OF_1=1\)、\(OF_2=3\)，所以由椭圆定义 \(2a=OF_1+OF_2=4\)，得 \(a=2\). 因而离心率
\(e=\frac ca=\frac12\)，选 C.
\end{solution}

\begin{problem}
某赛季足球比赛的计分规则是：胜一场，得 \(3\) 分；平一场，得 \(1\) 分；负一场，得 \(0\) 分. 一球队打完 \(15\) 场，积 \(33\) 分. 若不考虑顺序，该队胜，负，平的可能情况共有（\quad）

\choices
    {\(3\) 种}
    {\(4\) 种}
    {\(5\) 种}
    {\(6\) 种}
\end{problem}

\begin{answer}
A
\end{answer}

\begin{solution}
设胜、平、负的场数分别为 \(w\)，\(d\)，\(l\)，则
\(w+d+l=15\)，\(3w+d=33\).
消去 \(d\)，得 \(d=33-3w\)，\(l=2w-18\). 由 \(d\geq0\)、\(l\geq0\) 以及 \(w\geq0\)，得 \(9\leq w\leq11\). 因此 \(w\) 可取 \(9\)，\(10\)，\(11\)，每个 \(w\) 唯一确定 \(d\)，\(l\)，共有 \(3\) 种，选 A.
\end{solution}

\begin{problem}
设坐标原点为 \(O\)，抛物线 \(y^2 = 2x\) 与过焦点的直线交于 \(A\)，\(B\) 两点，则 \(\overrightarrow{OA} \cdot \overrightarrow{OB} =\)（\quad）

\choices
    {\(\dfrac{3}{4}\)}
    {\(-\dfrac{3}{4}\)}
    {\(3\)}
    {\(-3\)}
\end{problem}

\begin{answer}
B
\end{answer}

\begin{solution}
将抛物线上的点参数化为 \(X(t)=\left(\frac{t^2}{2},t\right)\). 设 \(A=X(t_1)\)、\(B=X(t_2)\). 抛物线上参数为 \(t_1\)，\(t_2\) 的两点所连弦的方程为
\((t_1+t_2)y=2x+t_1t_2\). 抛物线 \(y^2=2x\) 的焦点为 \(F\left(\frac12,0\right)\)，代入弦方程得 \(t_1t_2=-1\). 因此 \(\overrightarrow{OA}\cdot\overrightarrow{OB}=\frac{t_1^2t_2^2}{4}+t_1t_2=\frac14-1=-\frac34\)，故选 B.
\end{solution}

\begin{problem}
一间民房的屋顶有如图三种不同的盖法：
\begin{circlelist}
\item 单向倾斜；
\item 双向倾斜；
\item 四向倾斜.
\end{circlelist}
记三种盖法屋顶面积分别为 \(P_1\)，\(P_2\)，\(P_3\). 若屋顶斜面与水平面所成的角都是 \(\alpha\)，则（\quad）

\bitmapfigure[max width=0.60\linewidth]{img/2001/new_curriculum_liberal/q11_fig1.png}

\choices
    {\(P_3 > P_2 > P_1\)}
    {\(P_3 > P_2 = P_1\)}
    {\(P_3 = P_2 > P_1\)}
    {\(P_3 = P_2 = P_1\)}
\end{problem}

\begin{answer}
D
\end{answer}

\begin{solution}
设三种屋顶在水平面上的公共投影面积为 \(S_0\). 对于任一斜面，若斜面与水平面的夹角为 \(\alpha\)，则其水平投影面积等于斜面面积乘 \(\cos\alpha\). 三种盖法的各斜面投影共同恰好覆盖同一个水平屋顶，所以
\(P_1\cos\alpha=P_2\cos\alpha=P_3\cos\alpha=S_0\). 由于屋顶不是竖直面，\(\cos\alpha>0\)，故 \(P_1=P_2=P_3\)，选 D.
\end{solution}

\begin{problem}
如图，小圆圈表示网络的结点，结点之间的连线表示它们有网线相联. 连线标注的数字表示该段网线单位时间内可以通过的最大信息量. 现从结点 \(A\) 向结点 \(B\) 传递信息，信息可以分开沿不同的路线同时传递. 则单位时间内传递的最大信息量为（\quad）

\bitmapfigure[max width=0.42\linewidth]{img/2001/new_curriculum_liberal/q12_fig1.png}

\choices
    {\(26\)}
    {\(24\)}
    {\(20\)}
    {\(19\)}
\end{problem}

\begin{answer}
D
\end{answer}

\begin{solution}
按图在从 \(B\) 出发的四条支路上分别安排 \(3\)，\(4\)，\(6\)，\(6\) 单位的信息. 前两条支路经过的中间网线容量分别为 \(5\)，\(6\)，后两条支路经过的中间网线容量分别为 \(6\)，\(8\)，所以这些安排均不超过相应网线的容量. 前两条支路在右侧共用的容量为 \(12\) 的网线通过 \(3+4=7\) 单位，后两条支路在右侧共用的另一条容量为 \(12\) 的网线通过 \(6+6=12\) 单位，因而上述安排确实可以同时传递 \(3+4+6+6=19\) 单位的信息.

另一方面，割去下列四条网线后，\(B\) 与 \(A\) 不再连通：从 \(B\) 出发的最上方容量为 \(3\) 的网线、从 \(B\) 出发的第二条容量为 \(4\) 的网线、第三条分支通向右侧下方结点的容量为 \(6\) 的网线，以及从 \(B\) 出发的最下方容量为 \(6\) 的网线. 这个割的容量为 \(3+4+6+6=19\)，所以任何传输方案的总量都不超过 \(19\). 因此最大信息量为 \(19\)，选 D.
\end{solution}

\section{填空题}
\begin{problem}
定义在 \(\R\) 上的函数 \(f(x)=\sin x + \sqrt{3}\cos x\) 的最大值是\fillinblank{}.
\end{problem}

\begin{answer}
2
\end{answer}

\begin{solution}
由 \(\sin x+\sqrt3\cos x=2\left(\frac12\sin x+\frac{\sqrt3}{2}\cos x\right)=2\sin\left(x+\frac\pi3\right)\)，且正弦函数的最大值为 \(1\)，所以函数最大值为 \(2\).
\end{solution}

\begin{problem}
一个工厂有若干个车间，今采用分层抽样方法从全厂某天的 \(2048\) 件产品中抽取一个容量为 \(128\) 的样本进行质量检查. 若一车间这一天生产 \(256\) 件产品，则从该车间抽取的产品件数为\fillinblank{}.
\end{problem}

\begin{answer}
16
\end{answer}

\begin{solution}
分层抽样时，各层抽取数与该层产品数成比例. 抽样比例为
\(\frac{128}{2048}=\frac1{16}\)，所以第一车间抽取
\(256\times\frac1{16}=16\) 件.
\end{solution}

\begin{problem}
在空间中，
\begin{circlelist}
\item 若四点不共面，则这四点中任何三点都不共线.
\item 若两条直线没有公共点，则这两条直线是异面直线.
\end{circlelist}

以上两个命题中，逆命题为真命题的是\fillinblank{}.
\end{problem}

\begin{answer}
\(\circled{2}\)
\end{answer}

\begin{solution}
命题 \(\circled{1}\) 的逆命题是“若四点中任何三点都不共线，则这四点不共面”. 取同一平面内一个一般四边形的四个顶点，任何三点都不共线，但四点共面，所以逆命题 \(\circled{1}\) 为假.

命题 \(\circled{2}\) 的逆命题是“若两条直线是异面直线，则两条直线没有公共点”. 异面直线的定义正是不同在任何一个平面内且不相交的两条直线，所以该逆命题为真. 因此应填 \(\circled{2}\).
\end{solution}

\begin{problem}
设 \(\{a_n\}\) 是公比为 \(q\) 的等比数列，\(S_n\) 是它的前 \(n\) 项和. 若 \(|S_n|\) 是等差数列，则 \(q =\fillinblank{}\).
\end{problem}

\begin{answer}
1
\end{answer}

\begin{solution}
按本题所用等比数列的定义，首项 \(a_1\ne0\)，公比 \(q\ne0\). 设 \(c=|a_1|>0\). 由 \(S_1=a_1\)、\(S_2=a_1(1+q)\)、\(S_3=a_1(1+q+q^2)\)，且 \(q^2+q+1>0\)，得
\(|S_1|=c\)，\(|S_2|=c|1+q|\)，\(|S_3|=c(1+q+q^2)\).
因 \(|S_n|\) 为等差数列，前三项满足
\(2|S_2|=|S_1|+|S_3|\)，
即 \(2|1+q|=q^2+q+2\).

若 \(q\geq-1\)，则左边为 \(2q+2\)，从而 \(q^2-q=0\)，得 \(q=0\) 或 \(q=1\). 因为 \(q\ne0\)，故此时 \(q=1\). 若 \(q<-1\)，则上式化为 \(-2q-2=q^2+q+2\)，即 \(q^2+3q+4=0\)，其判别式 \(9-16<0\)，无实数解. 因此 \(q=1\).
\end{solution}

\section{解答题}
\begin{problem}
已知等差数列前三项为 \(a\)，\(4\)，\(3a\)，前 \(n\) 项的和为 \(S_n\)，\(S_k = 2550\).

\begin{enumerate}
\item 求 \(a\) 及 \(k\) 的值；
\item 求 \(\displaystyle\lim_{n\to\infty}\left(\frac{1}{S_1}+\frac{1}{S_2}+\cdots+\frac{1}{S_n}\right)\).
\end{enumerate}
\end{problem}

\begin{solution}
\begin{enumerate}
\item 由等差数列性质，\(a+3a=2\times4\)，得 \(a=2\)，公差 \(d=2\). 因而 \(S_k=2k+\frac{k(k-1)}{2}\times2=k^2+k=2550\)，解得 \(k=50\)（舍去 \(k=-51\)）.

\item \(S_n=n(n+1)\)，所以
\(\displaystyle\sum_{i=1}^{n}\frac1{S_i}=\frac1{1\times2}+\frac1{2\times3}+\cdots+\frac1{n(n+1)}=\left(1-\frac12\right)+\left(\frac12-\frac13\right)+\cdots+\left(\frac1n-\frac1{n+1}\right)=1-\frac1{n+1}\).
因此极限为 \(1\).
\end{enumerate}
\end{solution}

\begin{problem}
设计一幅宣传画，要求画面面积为 \(4840\mathrm{~cm}^2\)，画面的宽与高的比为 \(\lambda\)（\(\lambda < 1\)），画面的上，下各有 \(8\mathrm{~cm}\) 空白，左，右各有 \(5\mathrm{~cm}\) 空白. 怎样确定画面的高与宽尺寸，能使宣传画所用纸张面积最小？

\begin{solution}
由宽高比及画面面积可知 \(0<\lambda<1\). 设画面高为 \(x\mathrm{~cm}\)，则宽为 \(\lambda x\mathrm{~cm}\)，且 \(\lambda x^2=4840\). 纸张面积 \(S=(x+16)(\lambda x+10)=4840+(16\lambda+10)x+160\).
由 \(x=22\sqrt{10}/\sqrt{\lambda}\)，得
\[
S=5000+44\sqrt{10}\left(8\sqrt{\lambda}+\frac5{\sqrt{\lambda}}\right)
\]
由均值不等式，
\(8\sqrt{\lambda}+\dfrac5{\sqrt{\lambda}}\geq2\sqrt{40}\)，等号成立当且仅当
\(8\sqrt{\lambda}=\dfrac5{\sqrt{\lambda}}\). 因此 \(S\) 在 \(\lambda=5/8\) 时取最小值. 此时 \(x=88\)，画面宽为 \(\lambda x=55\)，故画面高、宽分别为 \(88\mathrm{~cm}\)、\(55\mathrm{~cm}\).
\end{solution}
\end{problem}

\begin{problem}
如图，用 \(A\)，\(B\)，\(C\) 三类不同的元件连接成两个系统 \(N_1\)，\(N_2\). 当元件 \(A\)，\(B\)，\(C\) 都正常工作时，系统 \(N_1\) 正常工作；当元件 \(A\) 正常工作且元件 \(B\)，\(C\) 至少有一个正常工作时，系统 \(N_2\) 正常工作. 已知元件 \(A\)，\(B\)，\(C\) 正常工作的概率依次为 \(0.80\)，\(0.90\)，\(0.90\).

\begin{enumerate}
\item 求系统 \(N_1\) 正常工作的概率 \(P_1\)；
\item 求系统 \(N_2\) 正常工作的概率 \(P_2\).
\end{enumerate}

\bitmapfigure[max width=0.55\linewidth]{img/2001/new_curriculum_liberal/q19_fig1.png}
\end{problem}

\begin{solution}
\begin{enumerate}
\item 由相互独立性，\(P_1=0.80\times0.90\times0.90=0.648\)，故系统 \(N_1\) 正常工作的概率为 \(0.648\).

\item \(P_2=0.80\left[1-(1-0.90)(1-0.90)\right]=0.792\)，故系统 \(N_2\) 正常工作的概率为 \(0.792\).
\end{enumerate}
\end{solution}

\section{选做题：解答题}

\begin{problem}
如图，以正四棱锥 \(V-ABCD\) 底面中心 \(O\) 为坐标原点建立空间直角坐标系 \(O-xyz\)，其中 \(Ox \parallel BC\)，\(Oy \parallel AB\). \(E\) 为 \(VC\) 中点，正四棱锥底面边长为 \(2a\)，高为 \(h_0\).

\begin{enumerate}
\item 求 \(\cos\left\langle \overrightarrow{BE},\overrightarrow{DE} \right\rangle\)；
\item 记面 \(BCV\) 为 \(\alpha\)，面 \(DCV\) 为 \(\beta\)，若 \(\angle BED\) 是二面角 \(\alpha-VC-\beta\) 的平面角，求 \(\cos\angle BED\) 的值.
\end{enumerate}

\bitmapfigure[max width=0.62\linewidth]{img/2001/new_curriculum_liberal/q20a_fig1.png}
\end{problem}

\begin{solution}
\begin{enumerate}
\item 取 \(B(a,a,0)\)，\(C(-a,a,0)\)，\(D(-a,-a,0)\)，\(V(0,0,h_0)\)，则 \(E=(-a/2,a/2,h_0/2)\). 因此 \(\overrightarrow{BE}=\left(-\frac{3a}{2},-\frac a2,\frac{h_0}{2}\right)\)，\(\overrightarrow{DE}=\left(\frac a2,\frac{3a}{2},\frac{h_0}{2}\right)\).
\[
\overrightarrow{BE}\cdot\overrightarrow{DE}=-\frac{3a^2}{4}-\frac{3a^2}{4}+\frac{h_0^2}{4}=\frac{h_0^2-6a^2}{4}
\]
又
\[
|\overrightarrow{BE}|=|\overrightarrow{DE}|=\frac12\sqrt{10a^2+h_0^2}
\]
从而
\[
\cos\left\langle\overrightarrow{BE},\overrightarrow{DE}\right\rangle=\frac{h_0^2-6a^2}{10a^2+h_0^2}
\]

\item 因为 \(\angle BED\) 是二面角 \(\alpha-VC-\beta\) 的平面角，所以 \(BE\)、\(DE\) 都垂直于棱 \(VC\). 又 \(\overrightarrow{CV}=(a,-a,h_0)\)，于是
\(\overrightarrow{BE}\cdot\overrightarrow{CV}=0\)，即
\(-\frac{3a^2}{2}+\frac{a^2}{2}+\frac{h_0^2}{2}=0\)，从而 \(h_0^2=2a^2\). 代回第（1）问的结果，得 \(\cos\angle BED=-\frac13\).
\end{enumerate}
\end{solution}

\begin{problem}
如图，在底面是直角梯形的四棱锥 \(S-ABCD\) 中，\(\angle ABC = 90^\circ\)，\(SA \perp\) 面 \(ABCD\)，\(SA = AB = BC = 1\)，\(AD = \frac{1}{2}\).

\begin{enumerate}
\item 求四棱锥 \(S-ABCD\) 的体积；
\item 求面 \(SCD\) 与面 \(SBA\) 所成的二面角的正切值.
\end{enumerate}

\bitmapfigure[float=inline,max width=0.48\linewidth]{img/2001/new_curriculum_liberal/q20b_fig1.png}
\end{problem}

\begin{solution}
\begin{enumerate}
\item 底面直角梯形 \(ABCD\) 的面积为 \(\frac12(BC+AD)AB=\frac12\left(1+\frac12\right)=\frac34\)，所以四棱锥体积为 \(\frac13\times1\times\frac34=\frac14\).

\item 延长 \(BA\)、\(CD\) 相交于 \(E\)，连接 \(SE\)，则 \(SE\) 是所求二面角的棱. 由 \(AD\parallel BC\)，三角形 \(EAD\) 与 \(EBC\) 相似，且 \(\frac{EA}{EB}=\frac{AD}{BC}=\frac12\). 又 \(EB=EA+AB\)，所以 \(EA=1\)，从而 \(EB=2\). 由于 \(SA\perp\) 面 \(ABCD\)，有 \(SE^2=SA^2+AE^2=2\)，\(SB^2=SA^2+AB^2=2\)，而 \(EB^2=4\)，故 \(SE^2+SB^2=EB^2\)，即 \(SE\perp SB\).

由于 \(SA\perp\) 面 \(ABCD\)，且 \(SA\subset\) 面 \(SEB\)，所以面 \(SEB\perp\) 面 \(EBC\)，交线为 \(EB\). 又 \(BC\perp EB\)，故 \(BC\perp\) 面 \(SEB\)，所以 \(SB\) 是 \(CS\) 在面 \(SEB\) 上的射影. 因为 \(SB\perp SE\)，由射影定理得 \(CS\perp SE\). 因而 \(SB\)、\(SC\) 分别是两侧面内垂直于棱 \(SE\) 的直线，\(\angle BSC\) 是所求二面角的平面角. 又 \(BC\perp SB\)，所以 \(\tan\angle BSC=\frac{BC}{SB}=\frac{1}{\sqrt2}=\frac{\sqrt2}{2}\).
\end{enumerate}
\end{solution}

\begin{problem}
已知函数 \(f(x) = x^3 - 3ax^2 + 2bx\) 在点 \(x = 1\) 处有极小值 \(-1\). 试确定 \(a\)，\(b\) 的值，并求出 \(f(x)\) 的单调区间.
\end{problem}

\begin{solution}
由 \(f(1)=-1\)，且 \(x=1\) 是可导函数的极小值点，得 \(f'(1)=0\)，所以
\[
\begin{cases}
1-3a+2b=-1,\\
3-6a+2b=0
\end{cases}
\]
两式相减得 \(-3a=-1\)，所以 \(a=\frac13\)；代入第一式得 \(b=-\frac12\). 此时 \(f(x)=x^3-x^2-x\)，且 \(f'(x)=3x^2-2x-1=(3x+1)(x-1)\).
故 \(f(x)\) 在 \((-\infty,-\frac13)\) 和 \((1,+\infty)\) 上为增函数，在 \((-\frac13,1)\) 上为减函数；特别地，\(f'(x)\) 在 \(x=1\) 左侧为负、右侧为正，故 \(x=1\) 确为极小值点.
\end{solution}

\begin{problem}
设 \(0 < \theta < \frac{\pi}{2}\)，曲线 \(x^2\sin\theta + y^2\cos\theta = 1\) 和 \(x^2\cos\theta - y^2\sin\theta = 1\) 有 \(4\) 个不同的交点.

\begin{enumerate}
\item 求 \(\theta\) 的取值范围；
\item 证明这 \(4\) 个交点共圆，并求圆半径的取值范围.
\end{enumerate}
\end{problem}

\begin{solution}
\begin{enumerate}
\item 令 \(u=x^2\)，\(v=y^2\)，联立两曲线方程，相加、相减得
\[
\begin{cases}
(\sin\theta+\cos\theta)u+(\cos\theta-\sin\theta)v=2,\\
(\sin\theta-\cos\theta)u+(\sin\theta+\cos\theta)v=0
\end{cases}
\]
解得 \(u=\sin\theta+\cos\theta\)，\(v=\cos\theta-\sin\theta\). 要有四个不同交点，必须且只需 \(x^2>0\)、\(y^2>0\)，此时 \(x\)、\(y\) 分别可取正、负两个值，共有四个组合. 对于 \(0<\theta<\frac\pi2\)，有 \(\sin\theta+\cos\theta>0\)，而
\(\cos\theta-\sin\theta>0\) 等价于 \(\theta<\frac\pi4\)，所以 \(0<\theta<\frac\pi4\).

\item 由上式，\(x^2+y^2=(\sin\theta+\cos\theta)+(\cos\theta-\sin\theta)=2\cos\theta\)，故四个交点都满足以原点为圆心、半径 \(r=\sqrt{2\cos\theta}\) 的圆方程，因而共圆. 当 \(0<\theta<\frac\pi4\) 时，\(\sqrt[4]{2}<r<\sqrt2\).
\end{enumerate}
\end{solution}
